\begin{abstract}
 %\jy{This paper mainly focus on detecting gender bias in multiple coreference systems: BerkeleyCoref, Stanford CoreNLP, UW end-to-end coreference system. We demonstarte that all of these coreference systems have gender bias issues through experiments. Furthermore, we propose a simple way, e.g., train a model on union data(original training data and the gender-reversed training data) to reduce gender bias.}

%Coreference resolution, a task consisting in clustering narrative phrases referring to the same entities, is considered fundamental in natural language processing. It requires reasoning that goes beyond grammar and syntax and has to take into account discourse, common-sense and semantic understanding. 

%Recent data-driven approaches have significantly advanced research in co-reference resolution but these systems carry the risk of inadvertently relying on societal stereotypes present in the data. 
We introduce a new benchmark, WinoBias, for coreference resolution focused on gender bias.
Our corpus contains Winograd-schema style sentences with entities corresponding to people referred by their occupation (e.g. the nurse, the doctor, the carpenter).
We demonstrate that a rule-based, a feature-rich, and a neural coreference system all link gendered pronouns to pro-stereotypical entities with higher accuracy than anti-stereotypical entities, by an average difference of 21.1 in F1 score.
%and is perfectly balanced in terms of gendered references to each entity (e.g. he, she, her, him).
%We show that current co-reference resolution systems perform significantly worse in the presence of gender atypical roles.
%(as defined using labor demographic statistics). % KW: can move this detail to intro or method section.
Finally, we demonstrate a data-augmentation approach that, in combination with existing word-embedding debiasing techniques, removes the bias demonstrated by these systems in WinoBias without significantly affecting their performance on existing coreference benchmark datasets.
Our dataset and code are avialable at \url{http://winobias.org}.

%Furthermore, we show that simple baselines that apply debiasing to the word embeddings used in some of these systems, along with gender anonymization, and data augmentation techniques, achieve promising results in alleviating this problem. %We systematically observe that given the same context, coreference systems make significantly different predictions for \texttt{female} and \texttt{male} entities. 
%For example, they cannot identify that a pronoun ``she'' refers to ``the leader'' as such a co-referent pair seldom appears on the training data.
%In this paper, we analyze the impact of gender bias in a coreference resolution corpus and three coreference resolution systems. We propose an approach to reduce gender stereotypes inadvertently encoded in these coreference systems. Our proposed method is easy to implement and does not need any modification to the training procedure. Experiments show that the proposed approach effectively reduce gender bias without affecting the accuracy of the original system. 
\end{abstract}